\section{Linear Cryptanalysis of ITUbee}

Linear cryptanalysis model is very similar to the differential model. Only difference is XOR and three-forked branch changes their role in the model. In linear model we add constraints for three-forked branch and pass the XOR operation. Since middle of the fiestel structure is symmetric and XOR and three-forked branch have same constraints, same model with differential cryptanalysis can be also used for linear cryptanalysis. Sketch of the one round MILP model of the linear cryptanalysis is given in Figure \ref{fig: ITUbeeLin}

\begin{figure}[!ht]
\centering
\resizebox{1\textwidth}{!}{%
\begin{circuitikz}
\tikzstyle{every node}=[font=\large]
\node [font=\large] at (2.5,16) {R1};
\draw (2.5,15.5) to[short] (2.5,14.5);
\draw  (2.5,14) circle (0.5cm);
\draw [short] (2.5,13.5) -- (2.5,11.5);
\draw [->, >=Stealth] (2.5,11.5) -- (4,11.5);
\draw  (4,12) rectangle  node {\large SLS} (5.25,11);
\draw [->, >=Stealth] (5.25,11.5) -- (6.5,11.5);
\draw [short] (2.5,14.5) -- (2.5,13.5);
\draw [short] (2,14) -- (3,14);
\draw  (7,11.5) circle (0.5cm);
\draw [->, >=Stealth] (7.5,11.5) -- (8.75,11.5);
\draw  (8.75,12) rectangle  node {\large L} (10,11);
\draw [->, >=Stealth] (10,11.5) -- (11.5,11.5);
\draw  (11.5,12) rectangle  node {\large SLS} (12.75,11);
\draw [->, >=Stealth] (12.75,11.5) -- (14.5,11.5);
\draw  (15,11.5) circle (0.5cm);
\node [font=\large] at (15,16) {R0};
\draw (15,15.5) to[short] (15,14.5);
\draw  (15,14) circle (0.5cm);
\node [font=\large] at (3,11.75) {};
\draw [short] (15,13.5) -- (15,12);
\draw [short] (6.5,11.5) -- (7.5,11.5);
\draw [short] (7,12) -- (7,11);
\draw [short] (14.5,11.5) -- (15.5,11.5);
\draw [short] (15,12) -- (15,11);
\draw [short] (14.5,14) -- (15.5,14);
\draw [short] (15,14.5) -- (15,13.5);
\draw [short] (15,11) -- (15,10.25);
\draw [short] (2.5,11.5) -- (2.5,10.25);
\draw [short] (2.5,10.25) -- (15,8.25);
\draw [short] (15,10.25) -- (2.5,8.25);
\node [font=\large] at (1.5,14) {KL};
\node [font=\large] at (16,14) {KR};
\node [font=\large] at (7,12.5) {KR};
\node [font=\large] at (3.25,12) {R1a};
\node [font=\large] at (5.75,12) {R1b};
\node [font=\large] at (8,12) {R1b};
\node [font=\large] at (10.75,12) {R1c};
\node [font=\large] at (13.75,12) {R0};
\node [font=\large] at (2,11) {R2};
\node [font=\large] at (2.5,8) {R0};
\node [font=\large] at (15,8) {R2};
\end{circuitikz}
}%
\caption{ITUbee MILP linear cryptanalysis sketch}
\label{fig: ITUbeeLin}
\end{figure}

Best linear bias for ITUbee S-box input-output mask is $2^{-4}$. In our MILP model we found that minimum number of active s-box for 3 consecutive round is 16. We use piling up lemma in \cite{matsui1993linear} for calculating 3 round probability. 3 round probability is $2^{15}*(2^{-4})^{16}=2^{-49}$. Linear cryptanalysis is known-plaintext attack. Number of required plaintexts for the three round attack is $(2^{-49})^{-2}= 2 ^{96}$ which is more than all possible plaintexts for the cipher. As a conclusion, we can say that 3 round of the ITUbee algorithm is secure against linear cryptanalysis.    